% ---------------------------------------------------------------------------
% Encoder of AnemoiModelEncProcDec (GraphTransformerForwardMapper), drawn at the
% block level of "Attention Is All You Need" Fig. 1, with every wire annotated
% by its array shape.
%
% The tikzpicture only -- \input it into a document whose preamble loads
% notation.tex and the tikz libraries listed in network_diagram.tex.
%
% Layout: the input assembly on the far left (field encoder, forcings), then three
% lanes flowing left to right --
%   DATA   (top)     x-space nodes:  psi -> Linear -> LayerNorm -> key, value
%   EDGE   (middle)  encoder edges:  a_e -> Linear -> eps
%   LATENT (bottom)  z-space nodes:  a_z -> Linear -> LayerNorm -> query, self
% converging at the graph attention block, after which the latent stream
% continues alone along the middle line.
%
% Every node is placed off the \row*/\col* registers below, so the figure is
% reflowed by editing numbers in one place rather than by moving nodes.
% ---------------------------------------------------------------------------

% --- lane (row) positions ---------------------------------------------------
\def\rowdata{0}          % data-space lane
\def\rowedge{-6.0}       % edge lane; also the latent stream after attention
\def\rowlat{-12.0}        % latent-space lane
\def\rowbypass{4.2}      % raw psi running over the top to the decoder
\def\rowconn{-9.0}       % graph connectivity entering the attention block
\def\rowselfskip{-16.0}  % self branch routed under the attention block
\def\rowdstskip{-18.6}   % pre-norm latent skip routed under everything
\def\rowffskip{-3.0}     % skip around the feed forward

% --- column positions -------------------------------------------------------
\def\colin{-14.0}   % raw inputs, before assembly
\def\colfe{-8.0}     % field encoder
\def\colA{0}        % inputs
\def\colcat{4.8}    % concatenation
\def\colB{7.7}      % embedding Linear
\def\colC{11.4}     % LayerNorm
\def\colD{15.4}     % q / k / v / edge projections
\def\colsplit{13.4} % where a LayerNorm output branches in two
\def\colE{21.6}     % graph attention
\def\colF{26.6}     % residual 1
\def\colG{29.4}     % projection
\def\colH{32.6}     % residual 2
\def\colI{35.6}     % LayerNorm + feed forward
\def\colJ{38.6}     % residual 3
\def\colK{40.8}     % output

\def\stack{1.15}    % vertical offset of paired sub-boxes

\begin{tikzpicture}[
  x=1cm, y=1cm, font=\large,
  blk/.style     = {rounded corners=3pt, draw=trapline, fill=trapfill,
                    line width=1.1pt, align=center, inner sep=5pt,
                    minimum height=1.2cm, minimum width=2.1cm},
  attnblk/.style = {blk, fill=attnfill},
  fieldblk/.style = {blk, draw=dataclr, fill=dataclr!8, line width=1.6pt},
  sum/.style     = {circle, draw=trapline, fill=white, line width=1.1pt,
                    inner sep=0pt, minimum size=0.72cm},
  term/.style    = {align=center, inner sep=2pt},
  wire/.style    = {-{Stealth[length=3.4mm,width=2.8mm]}, line width=1.1pt,
                    line cap=round, rounded corners=5pt},
  plain/.style   = {line width=1.1pt, line cap=round, rounded corners=5pt},
  shp/.style     = {font=\normalsize, inner sep=2pt, fill=white,
                    fill opacity=0.85, text opacity=1},
  tag/.style     = {font=\normalsize\itshape, inner sep=2pt},
  lane/.style    = {font=\small\bfseries, anchor=south west},
]

% =========================== lane labels ====================================
% anchored on the left edge of the input boxes so nothing overhangs the
% bounding box (the standalone crop would clip it)
\node[lane, text=edgeclr]   at ({\colA-1.8}, {\rowedge+0.55})        {edges};
\node[lane, text=latentclr] at ({\colA-1.8}, {\rowlat+0.55})         {latent};

% =========================== INPUT ASSEMBLY =================================
% Observations reach the network only through the field encoder; the forcings pass
% through. Static forcings are taken once, at t; dynamic ones also at the target time.
\node[blk, text width=3.4cm] (obsin) at (\colin, {\rowdata+3.2})
  {observations\\[1pt]{\small per field and group}};
\node[blk, text width=3.4cm] (dynin) at (\colin, {\rowdata+1.6})
  {dynamic forcings};
\node[blk, text width=3.4cm] (statin) at (\colin, {\rowdata+0.2})
  {static forcings};
\node[lane, text=dataclr] at (obsin.north west) {data};
\node[fieldblk, text width=3.8cm] (fenc) at (\colfe, {\rowdata+3.2})
  {\textbf{Field encoder}\\[1pt]{\small MLP per field, mean over\\[-1pt]groups with count $>0$}};
\node[blk, text width=3.4cm, minimum height=3.8cm] (uin) at (\colA, {\rowdata+1.7})
  {flatten time,\\concatenate};

\draw[wire, dataclr] (obsin) -- (fenc);
\draw[wire, dataclr] (fenc.east) -- node[shp, above]
  {\shapethree{\NT}{\Nx}{\Nfield\Nfc}} (fenc.east -| uin.west);
\draw[wire, dataclr] (dynin.east) -- node[shp, above, pos=0.72]
  {\shapethree{(\NT{+}\NTf)}{\Nx}{\Npass}} (dynin.east -| uin.west);
\draw[wire, dataclr] (statin.east) -- node[shp, above, pos=0.72]
  {\shapetwo{\Nx}{\Nstatic}\ at $t$} (statin.east -| uin.west);

% =========================== DATA LANE ======================================
\node[blk, text width=3.4cm] (axin) at (\colA, {\rowdata-\stack})
  {node attributes\\$\attrx$};
\node[sum] (cat)  at (\colcat, \rowdata) {\scriptsize$\|$};
\node[blk] (embx) at (\colB, \rowdata) {Linear};
\node[blk] (lnx)  at (\colC, \rowdata) {LayerNorm};
\node[blk] (link) at (\colD, {\rowdata+1.15}) {Linear};
\node[blk] (linv) at (\colD, {\rowdata-1.15}) {Linear};
\coordinate (xsplit) at (\colsplit, \rowdata);

\draw[wire, dataclr] (uin.east)  -- node[shp, above, pos=0.45]
  {\shapetwo{\Nx}{(\Npsi{-}\Nax)}} (cat);
\draw[wire, dataclr] (axin.east) -- node[shp, below, pos=0.45]
  {\shapetwo{\Nx}{\Nax}} (cat);
\draw[wire, dataclr] (cat) -- node[shp, above] {\shapetwo{\Nx}{\Npsi}}
                             node[tag, below, text=dataclr] {$\inputvec$} (embx);
\draw[wire, dataclr] (embx) -- node[shp, above] {\shapetwo{\Nx}{\Nphi}} (lnx);
\draw[plain, dataclr] (lnx.east) -- (xsplit);
\draw[wire, dataclr] (xsplit) |- (link.west);
\draw[wire, dataclr] (xsplit) |- (linv.west);

% the mapper returns its source input untouched; it is the decoder's dst input
\draw[wire, dataclr, densely dashed] (cat.north) |- (\colK, \rowbypass)
  node[shp, above, pos=0.62] {\shapetwo{\Nx}{\Npsi}}
  node[term, anchor=west, xshift=3pt, align=left, font=\small]
  {returned unchanged;\\destination input\\of the decoder};

% =========================== EDGE LANE ======================================
\node[blk, text width=3.4cm] (aein) at (\colA, \rowedge)
  {edge attributes\\$\attre$};
\node[blk] (line) at (\colD, \rowedge) {Linear};
\node[term, anchor=west, align=left, font=\small, text=edgeclr,
      fill=white, inner sep=2pt] (conn)
  at ({\colB-1.4}, \rowconn) {graph connectivity\ \ $2\!\times\!\NEE$};

\draw[wire, edgeclr] (aein) -- node[shp, above] {\shapetwo{\NEE}{\Nae}} (line);

% =========================== LATENT LANE ====================================
\node[blk, text width=3.4cm] (azin) at (\colA, \rowlat)
  {node attributes\\$\attrz$};
\node[blk] (embz) at (\colB, \rowlat) {Linear};
\node[blk] (lnz)  at (\colC, \rowlat) {LayerNorm};
\node[blk] (linq) at (\colD, {\rowlat+1.15}) {Linear};
\node[blk] (lins) at (\colD, {\rowlat-1.15}) {Linear};
\coordinate (zsplit) at (\colsplit, \rowlat);

\draw[wire, latentclr] (azin) -- node[shp, above] {\shapetwo{\Nz}{\Naz}} (embz);
\draw[wire, latentclr] (embz) -- node[shp, above] {\shapetwo{\Nz}{\Nphi}} (lnz);
\draw[plain, latentclr] (lnz.east) -- (zsplit);
\draw[wire, latentclr] (zsplit) |- (linq.west);
\draw[wire, latentclr] (zsplit) |- (lins.west);

% =========================== ATTENTION ======================================
\node[attnblk, minimum height=15.6cm, text width=5.4cm] (attn) at (\colE, \rowedge)
  {\textbf{Graph Multi-Head}\\[1pt]\textbf{Attention}\\[3pt]
   {\small $\NH$ heads $\times$ $\Nc$ channels}\\[10pt]
   {\large$\displaystyle
      \attn_{ij}=\operatorname*{softmax}_{i\in\nbhd(j)}
      \frac{\query_j\cdot(\keyv_i+\edgeemb_{ij})}{\sqrt{\Nc}}$}\\[9pt]
   {\large$\displaystyle
      \sum_{i\in\nbhd(j)}\attn_{ij}\,(\valuev_i+\edgeemb_{ij})$}};

\draw[wire, dataclr] (link.east) -- node[shp, above]
  {\shapethree{\Nx}{\NH}{\Nc}} node[tag, below, text=dataclr] {$\keyv$}
  (link.east -| attn.west);
\draw[wire, dataclr] (linv.east) -- node[tag, below, text=dataclr] {$\valuev$}
  (linv.east -| attn.west);
\draw[wire, edgeclr] (line.east) -- node[shp, above, pos=0.42]
  {\shapethree{\NEE}{\NH}{\Nc}} node[tag, below, pos=0.42, text=edgeclr] {$\edgeemb$}
  (line.east -| attn.west);
\draw[wire, latentclr] (linq.east) -- node[shp, above]
  {\shapethree{\Nz}{\NH}{\Nc}} node[tag, below, text=latentclr] {$\query$}
  (linq.east -| attn.west);
\draw[plain, edgeclr, densely dotted] (conn.east) -- (conn.east -| attn.west);

% =========================== POST-ATTENTION CHAIN ===========================
\node[sum] (res1) at (\colF, \rowedge) {\scriptsize$+$};
\node[blk] (proj) at (\colG, \rowedge) {Linear};
\node[sum] (res2) at (\colH, \rowedge) {\scriptsize$+$};
\node[blk, text width=2.4cm] (ff) at (\colI, \rowedge)
  {LayerNorm\\MLP};
\node[sum] (res3) at (\colJ, \rowedge) {\scriptsize$+$};
\node[term, anchor=west, align=left] (out) at ({\colK+0.2}, \rowedge)
  {$\latentvec$\\[2pt]{\small to processor}};

\draw[wire, latentclr] (attn.east) -- (res1);
\draw[wire, latentclr] (res1) -- (proj);
\node[tag, above=6pt of res1, text=trapline, align=center]
  {\small added \emph{before} the\\[-1pt]\small output projection};
\draw[wire, latentclr] (proj) -- node[shp, above] {\shapetwo{\Nz}{\Nphi}} (res2);
\draw[wire, latentclr] (res2) -- (ff);
\draw[wire, latentclr] (ff) -- (res3);
\draw[wire, latentclr] (res3) -- node[shp, above] {\shapetwo{\Nz}{\Nphi}} (out);


% residual 1 -- the self branch, added in attention space, ahead of the
% output projection (AIAYN draws that projection inside its attention box)
\draw[wire, latentclr] (lins.south) |- ({\colF}, \rowselfskip) -- (res1.south);
\node[shp] at ({(\colD+\colF)/2}, \rowselfskip) {\shapetwo{\Nz}{\Nphi}};
\node[tag, anchor=west, text=latentclr]
  at ({\colD+1.25}, {\rowlat-1.15}) {$\selfv$};

% residual 2 -- pre-norm skip taken from the latent embedding
\draw[wire, latentclr] (embz.south) |- ({\colH}, \rowdstskip) -- (res2.south);
\node[shp] at ({(\colC+\colH)/2}, \rowdstskip) {\shapetwo{\Nz}{\Nphi}};

% residual 3 -- pre-norm skip around the MLP
\draw[wire, latentclr] ({(\colH+\colI)/2}, \rowedge) |- ({\colJ}, \rowffskip)
  -- (res3.north);

\end{tikzpicture}
